\begin{remarks}[5.7.8]
\label{IV.2.s5.7.8-fr}
Sea $X$ un preesquema localmente noetheriano \emph{de dimensión} 1; equivale entonces
decir que $X$ es un preesquema de Cohen-Macaulay, que verifica $(S_1)$, o que
verifica una de las propiedades $(S_n)$ para $n\geq1$, en virtud de las
definiciones (5.7.1) y (5.7.2). En virtud de (5.7.5), para los preesquemas de
dimensión 1 equivale asimismo decir que $X$ es un preesquema de Cohen-Macaulay
o que carece de ciclo primo asociado inmerso. Por ejemplo, todo preesquema
localmente noetheriano \emph{reducido} de dimensión 1 es un preesquema de
Cohen-Macaulay.
\end{remarks}

\begin{proposition}[5.7.9]
\label{IV.2.s5.7.9-fr}
Sean $A$, $B$ dos anillos noetherianos, $\rho:A\to B$ un homomorfismo de
anillos y $M$ un $B$-módulo tal que $M_{[\rho]}$ sea un $A$-módulo de tipo
finito. Sea $\mathfrak p$ un ideal primo de $A$; los ideales primos de $B$
situados sobre $\mathfrak p$ y pertenecientes a $\operatorname{Supp}(M)$ son
finitos en número, y si $(\mathfrak q_i)_{1\leq i\leq n}$ es la familia de
estos ideales, se tiene
\begin{align}
\dim_{A_{\mathfrak p}}(M_{\mathfrak p})
&=\sup_i\dim_{B_{\mathfrak q_i}}(M_{\mathfrak q_i})
\tag{5.7.9.1}\label{IV.2.s5.7.9.1-fr}\\
\operatorname{prof}_{A_{\mathfrak p}}(M_{\mathfrak p})
&=\inf_i\operatorname{prof}_{B_{\mathfrak q_i}}(M_{\mathfrak q_i}).
\tag{5.7.9.2}\label{IV.2.s5.7.9.2-fr}
\end{align}
\end{proposition}

Si $\mathfrak b$ es el anulador de $M$, $B/\mathfrak b$ se identifica con un
sub-$A$-módulo de $\operatorname{End}_A(M_{[\rho]})$, y es por tanto de tipo
finito, puesto que $A$ es noetheriano; solo hay, pues, un número finito de
ideales primos de $B$ situados sobre $\mathfrak p$ y que contienen a
$\mathfrak b$, y estos son precisamente los que pertenecen a
$\operatorname{Supp}(M)$. Reemplazando $B$ por $B/\mathfrak b$, lo cual no
cambia los segundos miembros de (5.7.9.1) y (5.7.9.2) (por (\textbf 0, 16.1.9) y
(\textbf 0, 16.4.8)), se puede, pues, suponer que $B$ es una $A$-álgebra \emph{finita}.
Pongamos $S=A\setmin\mathfrak p$ y $B'=S^{-1}B$; $B'$ es un anillo noetheriano
semi-local cuyos ideales maximales son $S^{-1}\mathfrak q_i$
($1\leq i\leq n$), y como $M_{\mathfrak p}$ es un
$A_{\mathfrak p}$-módulo de tipo finito, se tiene
$\dim_{A_{\mathfrak p}}(M_{\mathfrak p})=\dim_{B'}(M_{\mathfrak p})$
(\textbf 0, 16.1.9). En estas condiciones, la relación (5.7.9.1) se convierte en un
caso particular de (\textbf 0, 16.1.7.4). Para demostrar la relación (5.7.9.2), la cuestión se reduce enseguida, como en (\textbf 0, 16.4.8), al caso en que
$\operatorname{prof}_{A_{\mathfrak p}}(M_{\mathfrak p})=0$, y el mismo
razonamiento que en (\textbf 0, 16.4.8) muestra que $M_{\mathfrak p}$ contiene un
sub-$B'$-módulo de \emph{longitud finita} y, por consiguiente, también un
sub-$B'$-módulo \emph{simple}; pero tal submódulo es necesariamente isomorfo al
cuerpo residual de uno de los $B_{\mathfrak q_i}$, por tanto existe al menos
un índice $i$ tal que
$\operatorname{prof}_{B_{\mathfrak q_i}}(M_{\mathfrak q_i})=0$
(\textbf 0, 16.4.6), lo que termina la demostración.
